% paper/sections/11_verification.tex
%
% §11  NS 方程式の物理的整合性検証 (Physical Consistency Verification)
%
% 検証の 3 層構造：
%   第10章 — コンポーネント単体（各演算子が設計精度を満たすか）
%   第11章 — 物理的整合性（組み上げた NS ソルバが物理法則を正しく記述するか）
%   第12章 — 二相流ベンチマーク（公表参照解との定量比較）
% ═══════════════════════════════════════════════════════════════════════════════
\section{NS 方程式の物理的整合性検証}
\label{sec:verification}%
\label{sec:grid_convergence}% backward-compat: referenced from §4 (curvature convergence → now in §10.1)

第~\ref{sec:numerical_integrity}章では，CCD/DCCD 微分演算・PPE ソルバ・TVD-RK3 等の
個々のコンポーネントが設計精度を有することを確認した．
本章では，これらのコンポーネントを組み合わせた
\textbf{ナビエ・ストークスソルバ全体}が物理法則を正しく記述するかを検証する．

\noindent\textbf{本章のスコープ：}
本章は\textbf{単相ナビエ・ストークスソルバ}の物理的整合性を対象とする．
すべてのテストケースは等密度（$\rho = \mathrm{const.}$）条件下で実施し，
Predictor--PPE--Corrector パイプラインの結合精度を検証する．
高密度比（$\rho_l/\rho_g \gg 1$）における界面横断時の整合性，
および二相流における質量保存性（密度加重）は，
GFM 統合後の物理ベンチマーク（第~\ref{sec:validation}章：静止液滴・毛細管波・気泡上昇・RT 不安定）にて定量的に検証する．

検証は以下の 4 段階で構成される：
\begin{enumerate}[noitemsep]
  \item 力学的バランス——静止流体における重力・表面張力の釣合（§\ref{sec:force_balance}）
  \item 保存則——運動エネルギー減衰と非圧縮条件の維持（§\ref{sec:ns_conservation}）
  \item 総合精度——NS 方程式の時間・空間収束次数（§\ref{sec:ns_time_accuracy}）
  \item 移流・圧力干渉——高レイノルズ数での対流不安定性に対する DCCD の効果（§\ref{sec:advection_pressure_coupling}）
\end{enumerate}

\noindent
各テストは理論解・解析解との定量比較により合否を判定する．
結果を踏まえた精度バジェットを §~\ref{sec:error_budget} に総括し，
第~\ref{sec:validation}章の二相流ベンチマークへと接続する．

\begin{table}[htbp]
\centering
\caption{物理的整合性の検証指標と許容基準}
\label{tab:verification}
\begin{tabular}{p{35mm}p{40mm}p{35mm}c}
\toprule
検証項目 & 理論値/解析解 & チェック方法 & 許容基準 \\
\midrule
静水圧バランス & $p = \rho g z$，$\bu = \bm{0}$ & $\|\bu\|_\infty$ & $< 10^{-12}$\\
Laplace 圧均衡 & $\Delta p = \sigma/R$ & 寄生流れ振幅 & $\Ord{h^2}$ 収束\\
運動エネルギー減衰 & $E_k(t) = E_k(0)\,e^{-4\nu t}$ & 相対誤差 & $< 10^{-4}$\\
非圧縮条件 & $\bnabla\cdot\bu \equiv 0$ & $\|\bnabla\cdot\bu\|_\infty$ & $< 10^{-10}$\\
NS 時間精度 & $\Ord{\Delta t^2}$（AB2+IPC） & $\Delta t$ 精緻化 & スロープ確認\\
NS 空間精度 & Kovasznay 解析解 & 格子精緻化 & スロープ確認\\
\bottomrule
\end{tabular}
\end{table}
